\section{Institutional Glossary and Statutory Framework}
\label{app:glossary}

This appendix compiles a comprehensive reference guide to Dutch labor law statutes, collective bargaining institutions, and statistical concepts utilized throughout the report.

\subsection*{Statutory baselines and collective bargaining institutions}

\begin{table}[H]
    \centering
    \caption{Glossary of statutory baselines and collective labor institutions}
    \label{tab:glossary_statutory}
    \small
    \begin{tabular}{lp{11.2cm}}
        \hline\hline
        Term / Statute & Legal Foundation and Substantive Economic Definition \\
        \hline
        \textbf{AVV} & \emph{Algemeen Verbindend Verklaring} (Act on Universally Binding Declarations of Collective Agreements, \emph{Wet AVV 1937}). A ministerial decree issued by the Minister of Social Affairs and Employment (SZW) extending sectoral agreement terms (\emph{bedrijfstak-cao}) to all employers and workers within that sector, eliminating wage undercutting. \\
        \textbf{TTW} & \emph{Tussentijdse Wijziging} (Interim Contract Amendment). A formal interim addendum agreed upon by signatory social partners during an active multi-year contract to modify specific provisions (such as wage grids or travel allowances) without terminating the underlying agreement. \\
        \textbf{BW (Book 7, Title 10)} & \emph{Burgerlijk Wetboek} (Dutch Civil Code, Articles 7:610--7:691). Governs statutory employment contracts. Key provisions operate as semi-mandatory law (\emph{driekwartdwingend recht}): individual contracts cannot undercut statutory floors, but collective agreements may legally deviate. \\
        \textbf{Ketenregeling} & Successive Fixed-Term Contract Chain Limits (Article~7:668a BW). Restricts consecutive fixed-term contracts before automatic conversion to permanent status (tightened to 3 contracts in 24 months under WWZ 2015; restored to 3 contracts in 36 months under WAB 2020; CAOs may expand up to 6 contracts in 48 months). \\
        \textbf{WML} & \emph{Wet minimumloon en minimumvakantiebijslag} (Statutory Minimum Wage and Minimum Holiday Allowance Act, 1968). Sets the mandatory adult wage floor and an 8\% holiday allowance. On 1 January 2024, fixed monthly minimums were abolished for a uniform hourly wage based on a 36-hour workweek. \\
        \textbf{WAZO} & \emph{Wet arbeid en zorg} (Work and Care Act). Governs statutory leaves, including 16 weeks maternity leave at 100\%, partner leave under WIEG (1 week at 100\% plus 5 weeks at 70\%), and paid parental leave under WBU (9 weeks at 70\% of earnings). CAOs frequently top up 70\% benefits to 100\%. \\
        \textbf{ATW} & \emph{Arbeidstijdenwet} (Working Hours Act). Establishes statutory maximum shifts (12 hours) and weekly hours (60 hours, averaging 48 hours over 16 weeks). CAOs enforce tighter weekly caps and negotiate unsocial/overtime premium rates. \\
        \textbf{O\&O fondsen} & \emph{Opleidings- en Ontwikkelingsfondsen} (Training and Development Funds). Bipartite sectoral foundation boards financed via collective payroll levies (0.2\%--0.8\%, extended via AVV) to fund industry-wide vocational training and eliminate worker poaching externalities. \\
        \textbf{PSA} & \emph{Psychosociale arbeidsbelasting} (Psychosocial Workload, Article~1(3) Working Conditions Act, \emph{Arbowet}). Workplace stress induced by sexual harassment, aggression, bullying, and extreme work pressure; employers must conduct risk inventories and designate confidential counsellors (\emph{vertrouwenspersonen}). \\
        \textbf{PMO / PAGO} & \emph{Periodiek Medisch Onderzoek} / \emph{Periodiek Arbeidsgezondheidskundig Onderzoek} (Article~18 Arbowet). Statutory preventive medical examinations offered confidentially to employees by certified occupational health doctors (\emph{arbodienst}). \\
        \textbf{Wet Bpf 2000} & \emph{Wet verplichte deelneming in een bedrijfstakpensioenfonds 2000} (Mandatory Industry Pension Funds Act). Authorizes ministerial decrees mandating participation in an industry-wide pension fund for all employers and workers in a sector. \\
        \hline\hline
    \end{tabular}
\end{table}

\subsection*{Wage structures, bargaining mechanisms, and statistical methodology}

\begin{table}[H]
    \centering
    \caption{Glossary of wage structures, bargaining conventions, and statistical indices}
    \label{tab:glossary_methodology}
    \small
    \begin{tabular}{lp{11.2cm}}
        \hline\hline
        Concept / Method & Substantive Definition and Methodological Operationalization \\
        \hline
        \textbf{Periodieken} & Pre-determined periodic (typically annual) salary step increments (\emph{ervaringsjaren} / \emph{treden}) awarded within a specific salary scale based on employee tenure or performance. \\
        \textbf{Pooled-z Score} & A standardization method centering and scaling provisions against the empirical multi-year corpus distribution ($z = (x - \mu)/\sigma$), enabling continuous tracking of secular trends without forcing annual cohort means to zero. \\
        \textbf{Winsorization} & Truncation of raw numeric provisions to the 1st and 99th percentiles ($[P_1, P_{99}]$) of the de-duplicated yardstick sample, followed by hard-clipping standardized z-scores to $[-3.0, +3.0]$ to contain outlier leverage. \\
        \textbf{KMO Measure} & Kaiser-Meyer-Olkin measure of sampling adequacy, assessing whether pairwise correlations between domains reflect shared latent factors suitable for factor analysis. \\
        \textbf{Statutory Elasticity} & The coefficient $\beta_0$ from CAO-fixed-effects regressions, measuring the contemporaneous standard-deviation adjustment in contractual terms per one-standard-deviation increase in statutory baselines. \\
        \hline\hline
    \end{tabular}
\end{table}

